% \begin{tikzpicture}

%     % ============================================================
%     % ÉCHELLES ET PARAMÈTRES
%     % ============================================================

%     \pgfmathsetmacro{\echelleY}{1.5}
%     \pgfmathsetmacro{\echelleX}{1.25}

%     \def\lam{4}
%     \def\amp{1}

%     % Deux instants :
%     % Le profil se déplace vers les x décroissants
%     \def\xUn{6.5}
%     \def\DeltaX{2.5}
%     \pgfmathsetmacro{\xDeux}{\xUn-\DeltaX}

%     \def\phi{-90}

%     % ============================================================
%     % CLIP
%     % ============================================================

%     \clip
%         (-1.5*\echelleX,-2.3*\echelleY)
%         rectangle
%         (11.0*\echelleX,2.3*\echelleY);

%     % ============================================================
%     % GRILLE
%     % ============================================================

%     \backgrowngrid
%         {-1}{-2}{10}{2}
%         {\echelleX}{\echelleY}
%         {\colorThree!50!\colorTwo}
%         {\colorTwo}
%         {0.5ex}

%     % ============================================================
%     % ÉQUATION
%     % ============================================================

%     \node[
%         fill=\colorThree!25!\colorTwo,
%         rounded corners=6pt,
%         inner sep=8pt,
%         font=\sffamily\bfseries\fontsize{18pt}{22pt}\selectfont,
%         text=\colorOne
%     ]
%     at (4.5*\echelleX,1.85*\echelleY)
%     {
%         $\bm{y(x,t)=g(x+ct)}$
%     };

%     % ============================================================
%     % PROFIL À t_1
%     % ============================================================

%     \draw[
%         line width=0.8ex,
%         color=\colorFour,
%         domain=-1:10,
%         samples=250,
%         smooth
%     ]
%     plot
%     (
%         \x*\echelleX,
%         {
%             \echelleY*\amp*
%             cos(360/\lam*\x+\phi)
%         }
%     );

%     % ============================================================
%     % PROFIL À t_2
%     %
%     % g(x+c t_2) = g((x+c t_1)+Delta x)
%     %
%     % Donc translation vers la GAUCHE
%     % ============================================================

%     \draw[
%         line width=0.8ex,
%         color=\colorSix,
%         domain=-1:10,
%         samples=250,
%         smooth
%     ]
%     plot
%     (
%         \x*\echelleX,
%         {
%             \echelleY*\amp*
%             cos(360/\lam*(\x+\DeltaX)+\phi)
%         }
%     );

%     % ============================================================
%     % POSITION DES POINTS IDENTIFIABLES
%     % ============================================================

%     \pgfmathsetmacro{\xPUn}{\xUn*\echelleX}
%     \pgfmathsetmacro{\xPDeux}{\xDeux*\echelleX}

%     \pgfmathsetmacro{\yP}{
%         \echelleY*\amp*
%         cos(360/\lam*\xUn+\phi)
%     }

%     % Point P à t_1
%     \fill[
%         color=\colorFour
%     ]
%     (\xPUn,\yP)
%     circle (0.14cm);

%     \node[
%         above right,
%         font=\sffamily\bfseries\fontsize{15pt}{18pt}\selectfont,
%         text=\colorFour
%     ]
%     at (\xPUn,\yP)
%     {$P_1$};

%     % Point P à t_2
%     \fill[
%         color=\colorSix
%     ]
%     (\xPDeux,\yP)
%     circle (0.14cm);

%     \node[
%         above left,
%         font=\sffamily\bfseries\fontsize{15pt}{18pt}\selectfont,
%         text=\colorSix
%     ]
%     at (\xPDeux,\yP)
%     {$P_2$};

%     % ============================================================
%     % PROJECTIONS
%     % ============================================================

%     \draw[
%         dashed,
%         line width=0.45ex,
%         color=\colorFour
%     ]
%     (\xPUn,\yP)
%     -- (\xPUn,0);

%     \draw[
%         dashed,
%         line width=0.45ex,
%         color=\colorSix
%     ]
%     (\xPDeux,\yP)
%     -- (\xPDeux,0);

%     % ============================================================
%     % DÉPLACEMENT DU POINT VERS LA GAUCHE
%     % ============================================================

%     \draw[
%         ->,
%         >=Latex,
%         line width=0.8ex,
%         color=\colorOne
%     ]
%     (\xPUn,0.85*\echelleY)
%     --
%     node[
%         midway,
%         above,
%         fill=\colorTwo,
%         inner sep=3pt,
%         font=\sffamily\bfseries\fontsize{14pt}{17pt}\selectfont,
%         text=\colorOne
%     ]
%     {
%         $\Delta x=-c(t_2-t_1)$
%     }
%     (\xPDeux,0.85*\echelleY);

%     % ============================================================
%     % LÉGENDE DES DEUX PROFILS
%     % ============================================================

%     \draw[
%         line width=0.8ex,
%         color=\colorFour
%     ]
%     (0.2*\echelleX,-1.55*\echelleY)
%     --
%     (1.0*\echelleX,-1.55*\echelleY)
%     node[
%         right,
%         text=\colorFour,
%         font=\sffamily\bfseries\fontsize{13pt}{16pt}\selectfont
%     ]
%     {$t=t_1$};

%     \draw[
%         line width=0.8ex,
%         color=\colorSix
%     ]
%     (3.2*\echelleX,-1.55*\echelleY)
%     --
%     (4.0*\echelleX,-1.55*\echelleY)
%     node[
%         right,
%         text=\colorSix,
%         font=\sffamily\bfseries\fontsize{13pt}{16pt}\selectfont
%     ]
%     {$t=t_2>t_1$};

%     % ============================================================
%     % AXE X
%     % ============================================================

%     \draw
%     (-1.3*\echelleX,0)
%     edge[
%         thick,
%         line width=0.8ex,
%         ->,
%         >=Triangle,
%         color=\colorOne
%     ]
%     node[
%         pos=1,
%         right,
%         font=\sffamily\bfseries\fontsize{16pt}{18pt}\selectfont,
%         text=\colorOne
%     ]
%     {$x$}
%     (10.8*\echelleX,0);

%     % ============================================================
%     % AXE Y
%     % ============================================================

%     \draw
%     (0,-2.0*\echelleY)
%     edge[
%         thick,
%         line width=0.8ex,
%         ->,
%         >=Triangle,
%         color=\colorOne
%     ]
%     node[
%         pos=1,
%         above,
%         font=\sffamily\bfseries\fontsize{16pt}{18pt}\selectfont,
%         text=\colorOne
%     ]
%     {$y$}
%     (0,2.1*\echelleY);

%     % ============================================================
%     % CONCLUSION GRAPHIQUE
%     % ============================================================

%     \node[
%         fill=\colorSix!15!\colorTwo,
%         rounded corners=5pt,
%         inner sep=7pt,
%         font=\sffamily\bfseries\fontsize{15pt}{18pt}\selectfont,
%         text=\colorSix
%     ]
%     at (7.0*\echelleX,-1.55*\echelleY)
%     {
%         propagation vers les $\bm{x}$ décroissants
%     };

% \end{tikzpicture}


% 3. Environnement unique et superposition directe
\begin{tikzpicture}
  % On dessine les deux éléments sur le MÊME repère absolu
  \clip[](-4, -4) rectangle (9 , 3) ;

  %\node[transform canvas={scale=1}] at (0,0) {
    \node[] at (0,0) {
    \begin{tikzpicture}

        

        % \draw[
        %     color=colorFour,
        %     line width=1.0ex,
        %     ->
        %     ] (-1,0) -- ({5*cos(45)},{2*sin(45)})
        %     node[pos = 0.5, above , font=\fontsize{20pt}{1pt}\selectfont\bfseries, text=colorFour  , rotate = 30 ]{$r$}
        % ;

        

        \newcommand{\maCourbe}{
            plot[
                smooth,
                tension=0.5
            ] coordinates {
                (-5,1)(-4,2)(-3,2)(-2,2)(-1,2.5)(0,2)(0.5,0.4)(1,-1)(1.5,0)(2,2)(2.5,1.5)(3,0.8)(4,1)(5,-1)(6,-1.5)(7,1)(8,2)(9,1)(10,1)(11,1)(15,1)
            }
        }

        \begin{scope}
        
            \clip[](-6.5, -2.5) rectangle (11 , 3) ;
            \draw[
                color=colorFour,
                line width=0.8ex
            ] \maCourbe;
        


            \draw[
                shift = {(-3, 0)},
                color=colorSix,
                line width=0.8ex
            ] \maCourbe;
        \end{scope}

        \draw
            (-1.3,0)edge[thick,line width=0.8ex,->,color=\colorOne]node[pos=0.9,below,font=\sffamily\bfseries\fontsize{16pt}{18pt}\selectfont,text=\colorOne]{$x$}(10.8,0)
        ;

        \draw
            (0,-2)edge[thick,line width=0.8ex,->,color=\colorOne]node[pos=0.85,left,font=\sffamily\bfseries\fontsize{16pt}{18pt}\selectfont,text=\colorOne]{$y$}(0,3.5)
        ;

        \draw[shift={(-1,2)}]
            (0,0)edge[thick,line width=0.8ex,<-,color=\colorThree]node[pos=0.5,above,font=\sffamily\bfseries\fontsize{14pt}{18pt}\selectfont,text=\colorThree]{$\Delta x = c ( t_1 - t_2 ) $}(3,0)
        ;

        \draw[shift={(0.5,-2)}]
            (0,0)edge[thick,line width=0.8ex,color=colorFour]node[pos=1,right,font=\sffamily\bfseries\fontsize{14pt}{18pt}\selectfont,text=\colorThree]{$y(x,t=t_1) $}(1,0)
        ;

        \draw[shift={(5.0,-2)}]
            (0,0)edge[thick,line width=0.8ex,color=colorSix]node[pos=1,right,font=\sffamily\bfseries\fontsize{14pt}{18pt}\selectfont,text=\colorSix]{$y(x,t=t_2) $}(1,0)
        ;


        \node[
            fill=\colorThree!25!\colorTwo,
            rounded corners=6pt,
            inner sep=8pt,
            font=\sffamily\bfseries\fontsize{13pt}{22pt}\selectfont,
            text=\colorOne
        ]
        at (7.5,3.05)
        {
            $\bm{y(x,t)=g(x+ct)}$
        };


        %\draw[thick, decorate, decoration={random steps, segment length=3pt, amplitude=0.5pt}] (0,0) to[bend left] (4,2);

        

    \end{tikzpicture}   
  };

  


  

  % Grille de repérage
  %\draw[line width=2.5ex , color = red] (-8, -5.5) rectangle (8 , 4.5) ;
  %\drawgrid{-7}{10}{-4}{4}{1}{1}
\end{tikzpicture}